\cleardoublepage{}
\begin{center}
    \bfseries \zihao{3} 摘~要
\end{center}

针对传统检索增强生成（Retrieval-Augmented Generation, RAG）模型在复杂多跳问答中面临的证据碎片化与生成幻觉问题，本文实现并分析了一套基于统一证据接口、置信度重检与后验验证的RAG框架原型。该框架将RAG视为一个兼顾正确率、成本与安全性的序列决策过程，重点考察基于重排序置信度的伪相关反馈策略（Confidence-based Pseudo-Relevance Feedback）与基于 CoVe 思想的后验验证机制。

在实现上，系统首先通过 \texttt{EvidenceUnit} 抽象将文本、表格与图谱片段映射到统一证据入口；随后以重排序得分作为置信度指标，用阈值控制二次检索触发；最后在生成后执行声明级核验，当发现声明与证据冲突或证据缺失时返回“无答案”（No-Answer）。

本文共组织了系统消融、基线对比、阈值敏感性、误差诊断与 Route A 可信基线补充实验。结果表明，纯文本基线在 HotpotQA 全量测试上的有效 F1 为 6.16，而异构序列化会引入显著格式噪声，直接拉低系统性能；Adaptive 控制组几乎不改变最终 F1，表明 Adaptive 模块仅作用于召回增强，其改进无法稳定传导至端到端生成质量；加入严格验证的全功能系统则引发灾难性的验证崩溃，使系统完全丧失可用性，拒答率达到 97.2\%，有效 F1 降至 0.06。进一步的 \texttt{deepseek-v4-flash} 真实 LLM 小消融显示，A2 仅将 F1 从 21.28 小幅提升至 22.17，而 A3/D 仍分别出现 50.0\%/57.0\% 的拒答率；Verification-Aware Retrieval (VAR) 在三次 50 样本 repeated-run 中取得 23.51 的平均 F1，并将拒答率稳定压低到 12.67\%，N=100 复核中也将 hard reject 的 48.0\% 拒答率降至 18.0\%/15.0\%。200 样本真实 verifier 阈值实验进一步表明，提高验证阈值会降低不安全接受率，但也会显著提高错误拒答率。与此同时，严格短答案提示反而导致 F1 退化到 2.01 左右，说明答案格式问题需要显式抽取模块而非单纯 prompt 压缩。本文从系统组合角度揭示了多模块 RAG 中的关键失效机制，并提出可复现的缓解路径，为复杂推理场景下的 RAG 设计提供了实证依据。

\cleardoublepage{}
\begin{center}
    \bfseries \zihao{3} Abstract
\end{center}

This thesis studies evidence gaps and hallucination in RAG systems. This work identifies and characterizes failure mechanisms in multi-module RAG systems. It implements and analyzes a RAG framework prototype based on a unified evidence interface, confidence-aware re-retrieval, and post-hoc verification. By formulating RAG as a sequential decision process that balances accuracy, cost, and safety, the framework focuses on confidence-based Pseudo-Relevance Feedback (PRF) and verification-constrained generation inspired by Chain-of-Verification (CoVe).

First, the system maps text, table, and graph snippets into a unified \texttt{EvidenceUnit} interface. Second, it uses reranking scores as a confidence signal to trigger secondary retrieval when the evidence chain is weak. Finally, it performs claim-level verification after generation and returns "No-Answer" when the answer is unsupported or contradicted by evidence.

Experiments show that the plain-text baseline reaches an effective F1 of 6.16 on HotpotQA, whereas heterogeneous serialization introduces significant format noise and reduces it to 4.59. Adaptive retrieval only acts on recall enhancement, and its gains do not reliably propagate to end-to-end generation quality. In contrast, strict verification causes a catastrophic verification collapse with a 97.2\% rejection rate and drives effective F1 down to 0.06. Additional \texttt{deepseek-v4-flash} follow-up experiments confirm the same trend under real LLM generation and verification: A2 only improves F1 from 21.28 to 22.17, while A3 and D still cause 50.0\% and 57.0\% rejection rates. Verification-Aware Retrieval (VAR) reduces excessive rejection in both repeated 50-sample runs and a 100-sample follow-up. A 200-sample verifier threshold study confirms the safety-usability tradeoff between unsafe acceptance and false rejection, and a strict short-answer study further shows that prompt compression alone cannot solve answer-format errors. Overall, this thesis reveals key failure mechanisms in multi-module RAG systems and provides reproducible mitigation paths for RAG design in complex reasoning scenarios.
